\hypertarget{ej3}{}
\noindent \textbf{Ejercicio 3.} Usando reglas de equivalencia (conmutatividad, asociatividad, De Morgan, etc.) determinar si los siguientes pares de fórmulas son equivalencias. Indicar en cada paso qué regla se utilizó.

\begin{enumerate}[label=\alph*)]
    \item $(p \vee q) \wedge (p \vee r) \quad \text{y} \quad \neg p \rightarrow (q \wedge r)$
    \item $\neg(\neg p) \rightarrow (\neg(\neg p \wedge \neg q)) \quad \text{y} \quad p \rightarrow (p \vee q)$
    \item $((True \wedge p) \wedge (\neg p \vee False)) \rightarrow \neg(\neg p \vee q) \quad \text{y} \quad p \wedge \neg q$
    \item $(p \vee (\neg p \wedge q)) \quad \text{y} \quad \neg p \rightarrow q$
    \item $p \rightarrow (q \wedge \neg(q \rightarrow r)) \quad \text{y} \quad (\neg p \vee q) \wedge (\neg p \vee (q \wedge \neg r))$
\end{enumerate}

\vspace{0.2cm}
{\color{gray}\hrule}
\vspace{0.4cm}

\textbf{Resolución:}

\begin{enumerate}[label=\textbf{\alph*)}]
    
    \item \textbf{$(p \vee q) \wedge (p \vee r) \quad \text{y} \quad \neg p \rightarrow (q \wedge r)$} \\
    Partimos de la primera expresión y aplicamos la propiedad \textbf{Distributiva} a la inversa (sacando factor común $p$):
    \[ (p \vee q) \wedge (p \vee r) \equiv p \vee (q \wedge r) \]
    Luego, por la definición de \textbf{Implicación} ($x \rightarrow y \equiv \neg x \vee y$), sabemos que $p \equiv \neg(\neg p)$, entonces:
    \[ \neg(\neg p) \vee (q \wedge r) \equiv \neg p \rightarrow (q \wedge r) \]
    \underline{\textbf{Son equivalentes.}}

    \vspace{0.3cm}
    \item \textbf{$\neg(\neg p) \rightarrow (\neg(\neg p \wedge \neg q)) \quad \text{y} \quad p \rightarrow (p \vee q)$} \\
    Trabajamos el lado izquierdo. Primero aplicamos \textbf{Doble Negación} en el antecedente: $\neg(\neg p) \equiv p$. \\
    Luego aplicamos las \textbf{Leyes de De Morgan} y \textbf{Doble Negación} en el consecuente: 
    \[ \neg(\neg p \wedge \neg q) \equiv \neg(\neg p) \vee \neg(\neg q) \equiv p \vee q \]
    Reescribiendo la expresión completa nos queda exactamente el lado derecho:
    \[ p \rightarrow (p \vee q) \]
    \underline{\textbf{Son equivalentes.}}

    \vspace{0.3cm}
    \item \textbf{$((True \wedge p) \wedge (\neg p \vee False)) \rightarrow \neg(\neg p \vee q) \quad \text{y} \quad p \wedge \neg q$} \\
    Simplificamos el antecedente usando la propiedad de \textbf{Elemento Neutro}:
    \[ True \wedge p \equiv p \]
    \[ \neg p \vee False \equiv \neg p \]
    Nos queda la conjunción $(p \wedge \neg p)$, que por la ley de \textbf{Contradicción} es siempre $False$. \\
    La implicación queda: $False \rightarrow \neg(\neg p \vee q)$. \\
    Como un antecedente falso hace que cualquier implicación sea automáticamente verdadera, todo el lado izquierdo se reduce a $True$. \\
    Como $True \not\equiv p \wedge \neg q$ (ya que la derecha depende de los valores de $p$ y $q$), concluimos que: \\
    \underline{\textbf{No son equivalentes.}}

    \vspace{0.3cm}
    \item \textbf{$(p \vee (\neg p \wedge q)) \quad \text{y} \quad \neg p \rightarrow q$} \\
    Aplicamos propiedad \textbf{Distributiva} en el lado izquierdo:
    \[ p \vee (\neg p \wedge q) \equiv (p \vee \neg p) \wedge (p \vee q) \]
    Por ley del \textbf{Tercero Excluido}, $(p \vee \neg p) \equiv True$. \\
    Por \textbf{Elemento Neutro}, $True \wedge (p \vee q) \equiv p \vee q$. \\
    Ahora miramos el lado derecho y aplicamos la definición de \textbf{Implicación}:
    \[ \neg p \rightarrow q \equiv \neg(\neg p) \vee q \equiv p \vee q \]
    Ambas expresiones se reducen a lo mismo. \\
    \underline{\textbf{Son equivalentes.}}

    \vspace{0.3cm}
    \item \textbf{$p \rightarrow (q \wedge \neg(q \rightarrow r)) \quad \text{y} \quad (\neg p \vee q) \wedge (\neg p \vee (q \wedge \neg r))$} \\
    Trabajamos primero el lado izquierdo. Aplicamos definición de \textbf{Implicación} en el paréntesis más interno:
    \[ p \rightarrow (q \wedge \neg(\neg q \vee r)) \]
    Aplicamos \textbf{De Morgan} y \textbf{Doble Negación}:
    \[ p \rightarrow (q \wedge (q \wedge \neg r)) \]
    Por propiedad \textbf{Asociativa} e \textbf{Idempotencia} ($q \wedge q \equiv q$):
    \[ p \rightarrow (q \wedge \neg r) \]
    Aplicamos definición de \textbf{Implicación} en la principal:
    \[ \neg p \vee (q \wedge \neg r) \]
    Ahora comparamos con el lado derecho. Si al lado derecho le aplicamos la propiedad \textbf{Distributiva} a la inversa (sacando $\neg p$ como factor común de la disyunción), obtenemos exactamente lo mismo:
    \[ (\neg p \vee q) \wedge (\neg p \vee (q \wedge \neg r)) \equiv \neg p \vee (q \wedge (q \wedge \neg r)) \equiv \neg p \vee (q \wedge \neg r) \]
    \underline{\textbf{Son equivalentes.}}

\end{enumerate}

\vspace{0.5cm}
\hrule
\vspace{0.5cm}